\documentclass[12pt]{article}
\usepackage[utf8]{inputenc}
\usepackage[spanish]{babel}
\usepackage{amsmath}
\usepackage{graphicx}
\usepackage{geometry}
\geometry{letterpaper, margin=1in}
\usepackage{enumitem}
\usepackage{titlesec}
\usepackage{fancyhdr}
\usepackage{setspace}

\usepackage[
backend=biber,
style=ieee,
sorting=ynt
]{biblatex}

\addbibresource{referencias.bib}

% Configuración de encabezados y pies de página
\pagestyle{fancy}
\fancyhf{}
\fancyhead[L]{\small Universidad Nacional de Altiplano\\Maestría en Informática}
\fancyhead[R]{\includegraphics[width=.9cm]{img/epgicono.png}}
\fancyfoot[C]{\thepage}

% Configuración del título
\titleformat{\section}
  {\normalfont\Large\bfseries}{\thesection}{1em}{}
\titleformat{\subsection}
  {\normalfont\large\bfseries}{\thesubsection}{1em}{}

% Título del documento
\title{\vspace{-2cm}
      \textbf{Actividad N° 00--}\\
      Denominación.}
\author{\textbf{Author}\\
        \small Universidad Nacional de Altiplano}
\date{}

\begin{document}
Actividad N° :...\\
\textbf{Maestrante : ...}
\doublespacing
\section*{Instrucciones Generales}
    \includegraphics[width=0.75\linewidth]{img/act1.png}



\newpage

\section{Variables}
Identifique y describa las variables principales que serán objeto de estudio en su investigación. Para cada variable, indique claramente su definición operativa y justifique su selección a partir de la revisión de la literatura:

\begin{enumerate}
    \item \textbf{Variable 1:} Descripción detallada y justificación con cita \cite{Kadhim2020165}.
    \item \textbf{Variable 2:} Descripción detallada y justificación con cita \cite{Paun20221}.
    \item \textbf{Variable 3:} Descripción detallada y justificación con cita \cite{Oliveira2020505}.
\end{enumerate}

\section{ Medida, Técnica e Instrumento a Utilizar}
Documente y justifique el tipo de medida, la técnica y el instrumento que utilizará para cada una de las variables de interés:

\begin{enumerate}
    \item \textbf{Tipo de Medida:} Definición precisa del tipo de medida (cualitativa o cuantitativa) y fundamentación teórica \cite{Xu2022139}.
    \item \textbf{Técnica:} Explicación detallada de la técnica de recolección de datos a utilizar, respaldada por literatura relevante \cite{Kadhim2020165}.
    \item \textbf{Instrumento:} Descripción del instrumento específico (cuestionario, entrevista, etc.) y su validación en estudios previos \cite{Awan20213}.
\end{enumerate}

\section{Técnicas e Instrumentos}
Documente la técnica e instrumentos a utilizar con referencias bibliográficas de alto impacto. Asegúrese de utilizar fuentes de SCOPUS, IEEE o ACM para garantizar la calidad y validez de la información citada.

\section{Operacionalización}

Construya la operacionalización de la variable con la información previamente estudiada.


\begin{itemize}
    \item \textbf{Tipo de Investigación:} Definición y justificación teórica \cite{Kadhim2020165}.
    \item \textbf{Método de Investigación:} Explicación detallada y ejemplos de uso en estudios previos \cite{Xu2022139}.
    \item \textbf{Diseño de Investigación:} Descripción y apoyo con referencias pertinentes \cite{Liang2020147}.
\end{itemize}

* Las referencias son de ejemplo.
\printbibliography

\end{document}
